% ============================================================
%  封面信息
% ============================================================
\school{计算机科学与技术学院}
\major{计算机科学与技术}
\student{2353105}{倪雨舒}
\thesistitle{Unix V6++ 虚拟存储器页表机制改造报告}{按页框管理与进程私有页表设计}
\thesistitleeng{Report on Page Table Redesign in Unix V6++}{Page-frame Allocation and Per-process Page Tables}
\thesisadvisor{操作系统课程设计}
\thesisdate{2026}{6}{20}

% 本课程设计不使用模板默认的“毕业设计（论文）”字样。
\renewcommand{\tjcovertext}{操作系统课程设计}
\renewcommand{\tjheadertext}{操作系统课程设计}

% 下面这些字段只在信息说明页或摘要页中使用。本报告不生成摘要页，
% 因此 main.tex 中没有调用 \MakeInfoPage、\MakeAbstract 或 \MakeAbstractEng。
\infotype{design}
\infoabstract{本报告围绕 Unix V6++ 中页表机制的改造展开，说明物理页框分配、进程私有页表、页目录切换以及 pageDemo 验证过程。}
\infodrawings{0}
\infowordcount{6000}
\infomaterials{
  \item 程序源代码；
  \item 运行截图。
}
